\documentclass[conference]{IEEEtran}
\IEEEoverridecommandlockouts
% The preceding line is only needed to identify funding in the first footnote. If that is unneeded, please comment it out.
\usepackage{cite}
\usepackage{amsmath,amssymb,amsfonts}
\usepackage{algorithmic}
\usepackage{graphicx}
\usepackage{textcomp}
\usepackage{xcolor}
\def\BibTeX{{\rm B\kern-.05em{\sc i\kern-.025em b}\kern-.08em
    T\kern-.1667em\lower.7ex\hbox{E}\kern-.125em X}}

% ---------------------------------------------------------------------------
% Paquetes adicionales (idioma, tablas, código, diagramas y enlaces)
% ---------------------------------------------------------------------------
\usepackage[utf8]{inputenc}
\usepackage[T1]{fontenc}
\usepackage[spanish,es-tabla,es-nodecimaldot,es-noshorthands,es-nolists]{babel}
\usepackage{booktabs}
\usepackage{array}
\usepackage{listings}
\usepackage{tikz}
\usetikzlibrary{arrows.meta,babel}
\usepackage{url}
\usepackage[hidelinks]{hyperref}

\addto\captionsspanish{\renewcommand{\figurename}{Fig.}}
\renewcommand{\IEEEkeywordsname}{Palabras clave}
\renewcommand{\lstlistingname}{Código}

% Tamaños de los símbolos matemáticos dentro de los bloques de código y de consola
\DeclareMathSizes{6}{6}{5}{5}
\DeclareMathSizes{6.4}{6.4}{5}{5}

% ---------------------------------------------------------------------------
% Estilos de código y de salida de consola
% ---------------------------------------------------------------------------
\definecolor{ckeyword}{RGB}{0,70,160}
\definecolor{ccomentario}{RGB}{40,120,40}
\definecolor{cfondo}{RGB}{24,24,28}
\definecolor{ctexto}{RGB}{220,220,220}
\definecolor{cemisor}{RGB}{95,210,215}
\definecolor{creceptor}{RGB}{215,120,215}
\definecolor{cvalido}{RGB}{110,215,110}
\definecolor{crechazo}{RGB}{240,95,95}
\definecolor{camarillo}{RGB}{235,200,80}
\definecolor{ctitulo}{RGB}{95,210,215}

% Tildes dentro de los bloques de código (funciona con pdfLaTeX y XeLaTeX)
\lstset{
  literate=
    {á}{{\'a}}1 {é}{{\'e}}1 {í}{{\'i}}1 {ó}{{\'o}}1 {ú}{{\'u}}1
    {Á}{{\'A}}1 {É}{{\'E}}1 {Í}{{\'I}}1 {Ó}{{\'O}}1 {Ú}{{\'U}}1
    {ñ}{{\~n}}1 {Ñ}{{\~N}}1 {ü}{{\"u}}1
}

\lstdefinestyle{csharp}{
  language=[Sharp]C,
  basicstyle=\fontfamily{lmtt}\fontsize{6.4}{7.3}\selectfont,
  keywordstyle=\color{ckeyword},
  commentstyle=\color{ccomentario}\itshape,
  morekeywords={var,byte},
  frame=tb, framerule=0.4pt, rulecolor=\color{black!40},
  columns=fullflexible, keepspaces=true, breaklines=true,
  showstringspaces=false, tabsize=4,
  aboveskip=4pt, belowskip=0pt,
  captionpos=b, abovecaptionskip=3pt, belowcaptionskip=0pt
}

% Salida de consola: <v>...</v> verde, <r>...</r> rojo, <a>...</a> amarillo, <t>...</t> títulos
\lstdefinestyle{consola}{
  basicstyle=\fontfamily{lmtt}\fontsize{6}{6.9}\selectfont\color{ctexto},
  basewidth={0.525em,0.45em},
  breaklines=true, breakindent=0pt,
  postbreak=\mbox{\textcolor{gray}{$\hookrightarrow$}\space},
  columns=fixed, keepspaces=true, showstringspaces=false,
  aboveskip=0pt, belowskip=0pt,
  emph={[1]EMISOR}, emphstyle={[1]\color{cemisor}},
  emph={[2]RECEPTOR}, emphstyle={[2]\color{creceptor}},
  moredelim=[is][\color{cvalido}]{<v>}{</v>},
  moredelim=[is][\color{crechazo}]{<r>}{</r>},
  moredelim=[is][\color{camarillo}]{<a>}{</a>},
  moredelim=[is][\color{ctitulo}]{<t>}{</t>}
}

% El fondo oscuro se pinta como un solo bloque (backgroundcolor de listings
% deja líneas claras entre renglones en algunos visores)
\newsavebox{\cajaconsola}

\begin{document}

\title{Firma digital con algoritmos propios de \textit{hash} y de llave privada transmitida por TCP en C\#}

\author{
\IEEEauthorblockN{Juan Mendez}
\IEEEauthorblockA{\textit{Facultad de Ingeniería} \\
\textit{Universidad de San Buenaventura}\\
Bogotá, Colombia}
\and
\IEEEauthorblockN{Natalia Noguera}
\IEEEauthorblockA{\textit{Facultad de Ingeniería} \\
\textit{Universidad de San Buenaventura}\\
Bogotá, Colombia}
\and
\IEEEauthorblockN{Karen Kemba}
\IEEEauthorblockA{\textit{Facultad de Ingeniería} \\
\textit{Universidad de San Buenaventura}\\
Bogotá, Colombia}
}

\maketitle

\begin{abstract}
Este documento presenta un sistema de firma digital en C\# sobre .NET que cifra el resumen de un mensaje con una llave privada y lo envía por TCP junto con la llave pública, para que el receptor verifique su autenticidad e integridad. El patrón adaptador permitió reemplazar las clases SHA256 y RSA de .NET por una función \textit{hash} propia de 256 bits y un cifrado propio de dos pasos basado en el diccionario de sustitución del grupo. La función propia obtuvo un efecto avalancha de 49,99\,\%, frente a 50,04\,\% de SHA-256, y el cifrado, una difusión de 49,53\,\%.
\end{abstract}

\begin{IEEEkeywords}
Firma digital, función \textit{hash}, RSA, patrón adaptador, TCP.
\end{IEEEkeywords}

% ===========================================================================
\section{Introducción}
% ===========================================================================

La firma digital permite verificar que un mensaje no fue modificado en tránsito (integridad) y que proviene de quien afirma enviarlo (autenticidad) \cite{menezes}. Se construye con una función \textit{hash}, que asigna a un mensaje de cualquier longitud un resumen fijo y debe ser determinista, irreversible y resistente a colisiones \cite{menezes,fips180}, con efecto avalancha: el cambio de un bit de la entrada altera cerca de la mitad de los bits de salida \cite{feistel}; y con criptografía de llave pública \cite{rsa}, en la que lo cifrado con la llave privada solo se descifra con la pública. El emisor cifra el resumen con su llave privada y el receptor lo compara con el que obtiene al descifrar la firma; si coinciden, el mensaje es íntegro y auténtico.

El laboratorio comprende dos etapas: el sistema con bibliotecas de C\#, organizado con el patrón adaptador \cite{gof}, y su sustitución por algoritmos propios. Se implementó como aplicación de consola de .NET~10 en lugar de Unity~3D; como los \textit{scripts} de Unity también se escriben en C\# \cite{unity}, la lógica criptográfica y de red puede trasladarse a Unity.

% ===========================================================================
\section{Desarrollo y evidencias}
\label{sec:desarrollo}
% ===========================================================================

\subsection{Arquitectura y patrón adaptador}

En la Fig.~\ref{fig:adaptador}, las interfaces objetivo \texttt{IHashAdapter} e \texttt{ICifradoAdapter} aíslan al cliente, \texttt{ServicioFirma}, de los algoritmos; cada adaptador traduce a esa interfaz una clase de .NET o propia, y \texttt{FabricaCripto} crea en el receptor el adaptador indicado por el nombre que viaja en el paquete. Así, el paso a los algoritmos propios no modificó \texttt{ServicioFirma} ni la red. En el Código~\ref{lst:verificar}, donde $p$ es el paquete recibido, el receptor obtiene los adaptadores, recalcula el resumen, descifra la firma con la llave pública y compara ambos resúmenes con \texttt{FixedTimeEquals}, cuyo tiempo depende de la longitud de los datos y no de su contenido \cite{msFixed}.

\begin{lstlisting}[style=csharp, caption={Verificación en ServicioFirma (simplificada).}, label={lst:verificar}]
var hash = FabricaCripto.CrearHash(p.AlgoritmoHash);
var cifrado = FabricaCripto
    .CrearCifradoVerificador(p.AlgoritmoCifrado);
byte[] calculado = hash.CalcularHash(UTF8(p.Mensaje));
byte[] recuperado = cifrado.DescifrarConLlavePublica(
    cifrado.TextoAFirma(p.Firma), p.LlavePublica);
bool valida = CryptographicOperations
    .FixedTimeEquals(calculado, recuperado);
\end{lstlisting}

\begin{figure}[!t]
\centering
\resizebox{0.9\columnwidth}{!}{%
\begin{tikzpicture}[font=\scriptsize,
  caja/.style={draw, align=center, text width=2.15cm, inner sep=2pt,
               minimum height=7mm, fill=#1},
  caja/.default=white,
  realiza/.style={-{Triangle[open, length=2mm, width=2mm]}, dashed},
  usa/.style={-{Stealth[length=1.8mm]}},
  adapta/.style={-{Stealth[length=1.8mm]}, densely dotted, thick}]
  % Cliente y fábrica
  \node[caja=yellow!15, text width=2.6cm] (srv) at (3.675,0) {\textbf{ServicioFirma}\\(cliente)};
  \node[caja=gray!12] (fab) at (7.35,0) {\textbf{FabricaCripto}\\(crea por nombre)};
  % Interfaces
  \node[caja=blue!10, text width=2.6cm] (ih) at (1.225,-0.95)
    {\guillemotleft interfaz\guillemotright\\\textbf{IHashAdapter}};
  \node[caja=blue!10, text width=2.6cm] (ic) at (6.125,-0.95)
    {\guillemotleft interfaz\guillemotright\\\textbf{ICifradoAdapter}};
  % Adaptadores
  \node[caja] (sha) at (0,-1.95) {\textbf{Sha256Adapter}};
  \node[caja] (hpa) at (2.45,-1.95) {\textbf{HashPropioAdapter}};
  \node[caja] (rsa) at (4.9,-1.95) {\textbf{RsaAdapter}};
  \node[caja] (cpa) at (7.35,-1.95) {\textbf{CifradoPropio-}\\\textbf{Adapter}};
  % Clases adaptadas
  \node[caja=orange!15] (net1) at (0,-2.95) {\textbf{SHA256}\\(.NET)};
  \node[caja=green!12] (pro1) at (2.45,-2.95) {\textbf{HashPropio}\\(propia)};
  \node[caja=orange!15] (net2) at (4.9,-2.95) {\textbf{RSA}, \textbf{BigInteger}\\(.NET)};
  \node[caja=green!12] (pro2) at (7.35,-2.95) {\textbf{CifradoPropio},\\\textbf{Diccionario} (propias)};
  % Relaciones
  \draw[usa] (srv) -- (fab);
  \draw[usa] (srv) -- (ih);
  \draw[usa] (srv) -- (ic);
  \draw[realiza] (sha.north) -- (ih.south);
  \draw[realiza] (hpa.north) -- (ih.south);
  \draw[realiza] (rsa.north) -- (ic.south);
  \draw[realiza] (cpa.north) -- (ic.south);
  \draw[adapta] (sha) -- (net1);
  \draw[adapta] (hpa) -- (pro1);
  \draw[adapta] (rsa) -- (net2);
  \draw[adapta] (cpa) -- (pro2);
\end{tikzpicture}}
\caption{Patrón adaptador. Las flechas discontinuas con triángulo indican implementación de la interfaz; las punteadas, la clase adaptada, y las continuas, uso. En naranja, clases de .NET; en verde, clases propias.}
\label{fig:adaptador}
\end{figure}

\begin{figure*}[!t]
\begin{lrbox}{\cajaconsola}
\begin{minipage}{\dimexpr\textwidth-8pt\relax}
\begin{lstlisting}[style=consola]
<t>  PUNTO 1 - LIBRERÍAS DE C#: SHA-256 + RSA-2048</t>
<a>--- Caso 1: mensaje legítimo (se espera FIRMA VÁLIDA) ---</a>
           Hash calculado del mensaje  : 1b21a2debd1e01a3db0f3a439eae82f341572e10f12b11f72ed52cc77abe2a03
           Hash descifrado de la firma : 1b21a2debd1e01a3db0f3a439eae82f341572e10f12b11f72ed52cc77abe2a03
<v>[RECEPTOR] FIRMA VÁLIDA. Los hashes coinciden: el mensaje es AUTÉNTICO (lo firmó el dueño de la llave privada) e ÍNTEGRO (no fue modificado).</v>
<a>--- Caso 2: un atacante cambia el texto en el camino (falla la INTEGRIDAD) ---</a>
           Hash calculado del mensaje  : d4ee5021ea57b0aac958c85671363d58da47f4914c4e69689af4ebf98e821d75
           Hash descifrado de la firma : 1b21a2debd1e01a3db0f3a439eae82f341572e10f12b11f72ed52cc77abe2a03
<r>[RECEPTOR] FIRMA RECHAZADA. Los hashes NO coinciden: el mensaje fue alterado en el camino o la firma no corresponde a esta llave pública.</r>
<a>--- Caso 3: un impostor firma con SU llave privada y se hace pasar por el emisor (falla la AUTENTICIDAD) ---</a>
           Hash descifrado de la firma : (no se pudo descifrar)
<r>[RECEPTOR] FIRMA RECHAZADA. No se pudo descifrar la firma con la llave pública: el bloque descifrado no tiene formato PKCS#1, la firma no se creó con la llave privada de esta llave pública. La firma NO es auténtica.</r>
<t>  PUNTO 2 - ALGORITMOS PROPIOS: HashPropio-256 + CifradoPropio (clave "gato")</t>
<a>--- Caso 1: mensaje legítimo (se espera FIRMA VÁLIDA) ---</a>
              Firma como se envía      : XD-mucho-vida-puffle-fife-memotenpiedad-puffle-bachata-fife-amigues-vida-encanta-puffle-fi... (419 caracteres)
           Hash calculado del mensaje  : f1813d91f9acec72f52715450f1b17162d18425247111c946f84e9602ef053aa
           Hash descifrado de la firma : f1813d91f9acec72f52715450f1b17162d18425247111c946f84e9602ef053aa
<v>[RECEPTOR] FIRMA VÁLIDA. Los hashes coinciden: el mensaje es AUTÉNTICO (lo firmó el dueño de la llave privada) e ÍNTEGRO (no fue modificado).</v>
<a>--- Caso 2: un atacante cambia el texto en el camino (falla la INTEGRIDAD) ---</a>
           Hash calculado del mensaje  : 6852fa8b42264f1ac2f498c3a7be45a3141f994f1280b7c37bf0137864f8b2f3
           Hash descifrado de la firma : f1813d91f9acec72f52715450f1b17162d18425247111c946f84e9602ef053aa
<r>[RECEPTOR] FIRMA RECHAZADA. Los hashes NO coinciden: el mensaje fue alterado en el camino o la firma no corresponde a esta llave pública.</r>
<a>--- Caso 3: un impostor firma con SU llave privada y se hace pasar por el emisor (falla la AUTENTICIDAD) ---</a>
           Hash descifrado de la firma : 1c595bf2586cde9b16d313725e99f6044d43b8a18160541ccd759083ff432ed4
<r>[RECEPTOR] FIRMA RECHAZADA. Los hashes NO coinciden: el mensaje fue alterado en el camino o la firma no corresponde a esta llave pública.</r>
\end{lstlisting}
\end{minipage}
\end{lrbox}
\setlength{\fboxsep}{4pt}%
\colorbox{cfondo}{\usebox{\cajaconsola}}
\caption{Extracto del modo de demostración: tres casos por etapa. Se omiten líneas intermedias; la salida completa se encuentra en el repositorio.}
\label{fig:demo}
\end{figure*}

\subsection{Comunicación por TCP}

El receptor escucha con \texttt{TcpListener} y el emisor se conecta con \texttt{TcpClient} \cite{msApi} para enviar un paquete JSON con el mensaje, la firma, la llave pública y los algoritmos. Como TCP transmite un flujo de bytes sin delimitar mensajes \cite{rfc9293}, cada paquete se precede de 4 bytes con su longitud, limitada a 1~MiB. Se prefirió TCP sobre UDP porque garantiza la entrega completa y ordenada \cite{rfc9293,rfc768}, sin la cual la verificación fallaría.

\subsection{Evidencias de funcionamiento}

La Fig.~\ref{fig:demo} muestra el modo de demostración con tres casos por etapa: el mensaje legítimo se acepta, el alterado en tránsito se rechaza por falta de integridad y el de un impostor, firmado con otra llave privada, por falta de autenticidad.

En la primera etapa, \texttt{SignHash} calcula $s = \mathrm{EM}^{d} \bmod n$, donde EM es el resumen con el relleno PKCS\#1~v1.5 \cite{rfc8017}; como \texttt{VerifyHash} solo devuelve un valor lógico, \texttt{RsaAdapter} calcula $s^{e} \bmod n$ con \texttt{BigInteger} para mostrar el resumen descifrado. En la segunda etapa, el adaptador del cifrado propio aplica el Código~\ref{lst:cifrado} (Sección~\ref{sec:justificacion}), donde $m$ es el resumen y $K$ la llave: un ciclo recorre los bytes hacia adelante, otro hacia atrás y el último los cambia de posición; descifrar aplica las inversas en orden contrario. La firma y la llave viajan como palabras del diccionario. También se verificó la combinación de HashPropio-256 con RSA y la ejecución en dos procesos independientes.

\begin{lstlisting}[style=csharp, caption={Cifrado en CifradoPropio (simplificado).}, label={lst:cifrado}]
byte ant = K[31];                          // ida
for (int i = 0; i < L; i++)
    ant = a[i] = Rotl((byte)((m[i] ^ K[i % 32]) + ant));
byte sig = K[0];                           // vuelta
for (int i = L - 1; i >= 0; i--)
    sig = c[i] = Rotl((byte)(a[i] + sig));
for (int i = 0; i < L; i++)                // paso 2
    f[(i * salto + inicio) % L] = c[i];
\end{lstlisting}

\subsection{Pruebas de los algoritmos propios}

La Tabla~\ref{tab:pruebas} resume las pruebas, ejecutadas con semillas fijas. En las de avalancha y difusión se invirtió un bit en 10\,000 entradas aleatorias y se contó cuántos bits del resultado cambiaron; el ideal es 128 de 256.

\begin{table}[!t]
\caption{Pruebas de los algoritmos propios}
\label{tab:pruebas}
\centering
\scriptsize
\begin{tabular}{@{}l>{\raggedright\arraybackslash}p{4.9cm}@{}}
\toprule
\textbf{Prueba} & \textbf{Resultado} \\
\midrule
Determinismo & Resumen idéntico en 3 ejecuciones \\
Avalancha (10\,000 mensajes) & 49,99\,\% (SHA-256: 50,04\,\%) \\
Distribución de bits & 48,62\,\%--51,64\,\% de unos por posición \\
Colisiones (200\,000 mensajes) & Ninguna \\
Reversibilidad (10\,000 casos) & 10\,000 de 10\,000 recuperados \\
Llave de otra palabra & No recupera el resumen \\
Llaves de ``gato'' y ``gatp'' & 132 de 256 bits distintos \\
Difusión del cifrado & 49,53\,\% (solo ida: 11,17\,\%) \\
RSA con \texttt{BigInteger} & Resumen original; \texttt{VerifyHash} acepta \\
\bottomrule
\end{tabular}
\end{table}

% ===========================================================================
\section{Justificación del diseño de los algoritmos}
\label{sec:justificacion}
% ===========================================================================

\subsection{Función \textit{hash} HashPropio-256}

El estado tiene ocho palabras de 32 bits (256 bits, como SHA-256), lo que hace intercambiables ambas funciones y fija en $2^{128}$ el costo genérico de una colisión \cite{menezes}. Las palabras iniciales son $h_k = \lfloor 2^{32}\,\mathrm{frac}(\sqrt{p_k}\,\pi)\rfloor$ para los primos 2 a 19, constantes de origen verificable análogas a las de SHA-256 \cite{fips180}. Las operaciones son módulo $2^{32}$.

\textit{Paso 1, absorción:} cada byte $b_i$ se combina con la palabra $h_j$, con $j = i \bmod 8$:
\begin{equation}
h_j \leftarrow (h_j \oplus b_i)\cdot P, \qquad h_{j+1} \leftarrow h_{j+1} + \mathrm{ROTL}_5(h_j).
\end{equation}
El XOR incorpora el byte sin perder información, el producto por el primo $P = \texttt{0x01000193}$ de FNV \cite{fnv} propaga los bits hacia las posiciones altas y la suma transmite el cambio a la palabra vecina. Al final, $h_0 \leftarrow h_0 \oplus |M|$.

\textit{Paso 2, mezcla:} se guarda una copia $h^{(0)}$ y se aplican cuatro vueltas de
\begin{align}
x &= \bigl(h_j + (h_{j-1} \oplus \mathrm{ROTL}_{11}(h_{j+3}))\bigr)\cdot Q, \\
h_j &\leftarrow x \oplus (x \gg 16),
\end{align}
con $Q = \texttt{0x9E3779B1} \approx 2^{32}/\varphi$ \cite{knuth}. El producto lleva los cambios a los bits altos y el desplazamiento los devuelve a los bajos; como cada palabra usa el valor ya actualizado de la anterior, un cambio alcanza las ocho palabras en una vuelta. El resumen es $H_j = h_j^{(4)} + h_j^{(0)}$.

La irreversibilidad se debe a la compresión a 256 bits, que descarta información; a la mezcla no lineal de sumas, XOR y rotaciones, y a la realimentación final, análoga a la de Davies--Meyer \cite{menezes}, que no se puede deshacer sin conocer $h^{(0)}$. Las pruebas estadísticas no sustituyen un criptoanálisis formal.

\subsection{Algoritmo de llave CifradoPropio}

Se limitó a dos pasos, uno de confusión y otro de difusión \cite{shannon}, para poder verificarlo a mano. La llave privada es una palabra $p$ y la pública es
\begin{equation}
K = \mathrm{HashPropio}\bigl(T(p)\bigr),
\end{equation}
donde $T$ traduce cada letra con el diccionario que el grupo desarrolló en \texttt{login.js}; por ejemplo, ``gato'' se traduce como ``muchovidapenguinde''. Así, una palabra corta produce 32 bytes de apariencia aleatoria de los que no se puede recuperar $p$.

\textit{Paso 1, sustitución encadenada de ida y vuelta:} con sumas módulo 256 y $\mathrm{ROTL}_3$ como rotación de 3 bits, la ida recorre el resumen $m$ de $i = 0$ a $L-1$ y la vuelta, de $L-1$ a $0$:
\begin{align}
a_i &= \mathrm{ROTL}_3\bigl((m_i \oplus K_{i \bmod 32}) + a_{i-1}\bigr), \\
c_i &= \mathrm{ROTL}_3(a_i + c_{i+1}),
\end{align}
con $a_{-1} = K_{31}$ y $c_L = K_0$. El XOR oculta cada byte con la llave y la suma con el vecino encadena los bytes, como en el modo CBC \cite{menezes}: la ida propaga un cambio hacia adelante y la vuelta, hacia atrás. La rotación evita que la diferencia avance como un desplazamiento constante que altera pocos bits; así, la difusión pasó de 11,17\,\% con solo la ida a 49,53\,\% (Tabla~\ref{tab:pruebas}). Se invierte con $a_i = \mathrm{ROTR}_3(c_i) - c_{i+1}$ y $m_i = \bigl(\mathrm{ROTR}_3(a_i) - a_{i-1}\bigr) \oplus K_{i \bmod 32}$.

\textit{Paso 2, transposición:} el byte de la posición $i$ pasa a
\begin{equation}
d(i) = (i\,s + b) \bmod L, \qquad \operatorname{mcd}(s, L) = 1,
\end{equation}
donde $L$ es el número de bytes, y $s$ y $b$ se derivan de $K$. La condición sobre $s$ hace biyectiva la permutación \cite{menezes}. Por ejemplo, ``Hola'' $= (72, 111, 108, 97)$ con ``gato'' da $(121, 213, 10, 152)$ tras el paso~1 y $(152, 10, 213, 121)$ tras el paso~2.

Para transmitir la firma y la llave, cada dígito hexadecimal se reemplaza por su palabra del diccionario, por ejemplo, \texttt{1}$\to$amigues y \texttt{a}$\to$vida; al \texttt{0}, ausente en el diccionario, se le asignó ``uwu''. Esta codificación no aporta seguridad, que debe depender solo de la llave \cite{kerckhoffs}.

\subsection{Limitaciones}

CifradoPropio es esencialmente simétrico: quien conoce $K$ podría firmar, mientras que en RSA obtener la llave privada exige factorizar $n$ \cite{rsa}. El diccionario no es inyectivo: ``ya'' y ``ja'' generan la misma llave. Además, un atacante podría sustituir la llave pública que viaja con el mensaje; el sistema muestra una huella para compararla por otro medio, problema que en la práctica resuelven los certificados X.509 \cite{rfc5280}.

% ===========================================================================
\section{Investigación de librerías}
\label{sec:librerias}
% ===========================================================================

La Tabla~\ref{tab:librerias} resume las bibliotecas empleadas, todas incluidas en .NET \cite{msApi}, por lo que no se requirieron dependencias externas. La clase \texttt{SHA256} rellena el mensaje hasta un múltiplo de 512 bits, incluida su longitud, y procesa cada bloque con 64 rondas de compresión sobre ocho palabras de 32 bits \cite{fips180}. En \texttt{RSA}, el módulo $n$ es el producto de dos primos grandes, el exponente público $e$ suele ser 65\,537 y el privado $d$ cumple $e\,d \equiv 1 \pmod{\varphi(n)}$ \cite{rsa}; su seguridad depende de la dificultad de factorizar $n$.

\begin{table}[!t]
\caption{Bibliotecas de .NET utilizadas}
\label{tab:librerias}
\centering
\scriptsize
\begin{tabular}{@{}>{\raggedright\arraybackslash}p{2.6cm}>{\raggedright\arraybackslash}p{5.8cm}@{}}
\toprule
\textbf{Clase} & \textbf{Uso y motivo de selección} \\
\midrule
\texttt{SHA256} & Resumen de la etapa 1; estándar FIPS 180-4 \cite{fips180} \\
\texttt{RSA} & Firma de 2048 bits con \texttt{SignHash}; la llave pública verifica, pero no firma \cite{rfc8017} \\
\texttt{BigInteger} & \texttt{ModPow}: $s^{e} \bmod n$ para mostrar el resumen \\
\texttt{Cryptographic\allowbreak Operations} & \texttt{FixedTimeEquals}: comparación en tiempo constante \\
\texttt{TcpListener}, \texttt{TcpClient} & Conexión confiable y ordenada \\
\texttt{JsonSerializer} & Serialización del paquete con validación de campos \\
\bottomrule
\end{tabular}
\end{table}

% ===========================================================================
\section{Conclusiones}
\label{sec:conclusiones}
% ===========================================================================

El sistema aceptó los mensajes legítimos y rechazó los alterados y los de un impostor con ambos pares de algoritmos, y el patrón adaptador permitió cambiar los algoritmos sin modificar el servicio de firma ni la red. HashPropio-256 igualó estadísticamente a SHA-256, aunque esto no demuestra su seguridad. La medición de la difusión mostró que encadenar bytes solo con sumas propaga poco los cambios y orientó el rediseño del cifrado. Por último, CifradoPropio cumplió la mecánica de la firma, pero el no repudio exige que la llave privada no pueda deducirse de la pública, como en RSA.

% ===========================================================================
\section{Repositorio}
\label{sec:repositorio}
% ===========================================================================

El código fuente se encuentra en \url{https://github.com/juanestebanmendez88-hash/lab3-firma-digital}, con las salidas completas de las ejecuciones (\texttt{informe/evidencias}), las instrucciones de uso (\texttt{README.md}) y un \texttt{.gitignore} de .NET que excluye los archivos generados (\texttt{bin/}, \texttt{obj/}).

% ===========================================================================
% Referencias
% ===========================================================================
\begin{thebibliography}{00}

\bibitem{menezes}
A.~J. Menezes, P.~C. van Oorschot, and S.~A. Vanstone, \emph{Handbook of Applied Cryptography}. Boca Raton, FL, USA: CRC Press, 1996.

\bibitem{fips180}
National Institute of Standards and Technology, ``Secure Hash Standard (SHS),'' FIPS PUB 180-4, Aug. 2015, doi: 10.6028/NIST.FIPS.180-4.

\bibitem{feistel}
H.~Feistel, ``Cryptography and computer privacy,'' \emph{Sci. Amer.}, vol.~228, no.~5, pp.~15--23, May 1973.

\bibitem{rsa}
R.~L. Rivest, A.~Shamir, and L.~Adleman, ``A method for obtaining digital signatures and public-key cryptosystems,'' \emph{Commun. ACM}, vol.~21, no.~2, pp.~120--126, Feb. 1978.

\bibitem{gof}
E.~Gamma, R.~Helm, R.~Johnson, and J.~Vlissides, \emph{Design Patterns: Elements of Reusable Object-Oriented Software}. Reading, MA, USA: Addison-Wesley, 1994.

\bibitem{unity}
Unity Technologies, ``Gameplay programming,'' Unity 6.6 User Manual. Accedido: 13 nov. 2026. [En línea]. Disponible: \url{https://docs.unity.com/en-us/engine/current/manual/scripting}

\bibitem{msFixed}
Microsoft, ``CryptographicOperations.FixedTimeEquals method,'' .NET API Reference. Accedido: 13 nov. 2026. [En línea]. Disponible: \url{https://learn.microsoft.com/dotnet/api/system.security.cryptography.cryptographicoperations.fixedtimeequals}

\bibitem{msApi}
Microsoft, ``.NET API browser,'' Microsoft Learn. Accedido: 13 nov. 2026. [En línea]. Disponible: \url{https://learn.microsoft.com/dotnet/api/}

\bibitem{rfc9293}
W.~Eddy, Ed., ``Transmission Control Protocol (TCP),'' IETF RFC 9293, Aug. 2022.

\bibitem{rfc768}
J.~Postel, ``User Datagram Protocol,'' IETF RFC 768, Aug. 1980.

\bibitem{rfc8017}
K.~Moriarty, B.~Kaliski, J.~Jonsson, and A.~Rusch, ``PKCS \#1: RSA Cryptography Specifications Version 2.2,'' IETF RFC 8017, Nov. 2016.

\bibitem{fnv}
L.~C. Noll, ``FNV hash.'' Accedido: 13 nov. 2026. [En línea]. Disponible: \url{http://www.isthe.com/chongo/tech/comp/fnv/index.html}

\bibitem{knuth}
D.~E. Knuth, \emph{The Art of Computer Programming, Vol. 3: Sorting and Searching}, 2nd~ed. Reading, MA, USA: Addison-Wesley, 1998.

\bibitem{shannon}
C.~E. Shannon, ``Communication theory of secrecy systems,'' \emph{Bell Syst. Tech. J.}, vol.~28, no.~4, pp.~656--715, Oct. 1949.

\bibitem{kerckhoffs}
A.~Kerckhoffs, ``La cryptographie militaire,'' \emph{J. Sci. Militaires}, vol.~9, pp.~5--38, Jan. 1883.

\bibitem{rfc5280}
D.~Cooper, S.~Santesson, S.~Farrell, S.~Boeyen, R.~Housley, and W.~Polk, ``Internet X.509 Public Key Infrastructure Certificate and Certificate Revocation List (CRL) Profile,'' IETF RFC 5280, May 2008.

\end{thebibliography}

\end{document}
